\chapter{Uma História de Ideias, Máquinas e Promessas}
\label{ch:historia}

Uma máquina ganha um jogo. Outra escreve um poema. Uma terceira decide quem entra numa fila de inspeção. Chamamos «inteligência artificial» às três, mas os problemas, mecanismos e consequências são diferentes. A história da área ajuda a perceber como chegámos a esta palavra tão elástica. Também funciona como vacina contra anúncios de que alguém resolveu, finalmente, «a inteligência».

\section{Antes de haver um nome}
Em 1950, Alan Turing discutiu a pergunta sobre máquinas que pensam através de um jogo de imitação~\citep{turing1950}. O artigo propõe uma forma de discutir comportamento observável e enfrenta objeções; não fornece um teste universal de consciência. Conseguir conversar convincentemente não determina, por si só, o mecanismo que produz a conversa ou a existência de experiência subjetiva.

A proposta de Dartmouth foi escrita em 1955 para um encontro no verão de 1956 por John McCarthy, Marvin Minsky, Nathaniel Rochester e Claude Shannon~\citep{dartmouth1955}. O nome da área ficou ligado a esse programa. Não significa que lógica, computação, estatística ou estudo do cérebro tenham começado naquela sala. Uma área científica ganha uma identidade pública em momentos que raramente coincidem com o início de todas as suas ideias.

É útil distinguir três perguntas históricas: podemos descrever raciocínios por regras? Podemos aprender regularidades a partir de exemplos? Podemos agir num ambiente para alcançar um objetivo? A IA moderna combina respostas a estas perguntas, em vez de ter substituído de uma vez por todas uma escola antiga por outra nova.

\section{Regras, procura e conhecimento explícito}
Num sistema simbólico, representamos elementos e relações de forma explícita. Uma regra pode dizer: se existe uma peça em falta e há talão válido, abrir investigação. Uma procura explora possibilidades, como percursos num mapa ou jogadas num tabuleiro. O programador não precisa de listar todas as sequências; define estados, movimentos e critérios.

O benefício é a possibilidade de inspecionar certas decisões. O custo surge quando o mundo resiste ao catálogo: descrições ambíguas, exceções, informação incompleta. A Lavandaria pode escrever cinquenta regras e descobrir uma reclamação que não encaixa em nenhuma. Isso não demonstra que regras são inúteis. Demonstra que precisamos de conhecer o domínio e de tratar casos fora das regras.

Sistemas baseados em conhecimento fizeram desta manutenção uma questão central. Se a política muda, quem altera as regras? Como se resolvem conflitos? A mesma pergunta regressa hoje quando alguém guarda uma política num prompt e se esquece de indicar a versão. A novidade da interface não elimina a antiguidade do problema.

\section{Aprender representações e aumentar a escala}
A aprendizagem automática desloca parte do trabalho: em vez de programar toda a associação entre entrada e saída, escolhemos dados, modelo e objetivo. Redes neuronais profundas aprendem representações em várias camadas. Transformers acrescentaram uma arquitetura de atenção que se tornou central em modelos de linguagem~\citep{vaswani2017}. Os capítulos seguintes abrem estes mecanismos; aqui interessa a mudança de método.

Mais dados e computação podem alargar capacidades. Não transformam automaticamente uma demonstração numa aplicação fiável. Há pelo menos três momentos diferentes: obter um resultado num conjunto de avaliação; fazê-lo funcionar sob condições operacionais; e demonstrar benefício para quem usa ou é afetado pelo sistema. Saltar do primeiro para o terceiro é uma forma bastante eficiente de escrever um comunicado e bastante ineficiente de fazer ciência.

Também não existe apenas uma arquitetura vencedora. Um sistema atual pode juntar pesquisa clássica, regras, rede neuronal, simulação e otimização. Um produto que conversa pode esconder uma função determinística para calcular preços. Isso costuma ser uma qualidade do desenho, não uma confissão de que a IA «não era verdadeira».

\section{ANI, AGI e ASI: palavras precisam de critérios}
\emph{ANI}, inteligência artificial estreita, designa habitualmente sistemas orientados a capacidades ou domínios delimitados. \emph{AGI}, inteligência artificial geral, refere uma ambição de competência ampla e transferência entre tarefas. \emph{ASI}, superinteligência artificial, é usada para capacidades superiores às humanas em domínios de grande amplitude. Não há aqui uma escala internacional única com três certificados para pendurar na parede.

Antes de aceitar uma alegação de generalidade, pergunta: que tarefas, com que ajuda, a que custo, com que aprendizagem adicional e com que taxa de falha? Um sistema pode dominar uma prova de matemática e falhar uma tarefa de navegação. Um resultado com ferramentas e tentativas ilimitadas não é equivalente ao de um modelo isolado numa única chamada. A comparação precisa de conservar essas condições.

Autonomia também admite graus: sugerir o próximo passo, escolher entre ferramentas autorizadas, executar uma sequência ou alterar os próprios planos sob limites. Nenhum desses comportamentos exige atribuir vontade humana ao sistema. A fronteira de ação deve ser descrita por permissões e condições de paragem, não por um adjetivo impressionante.

\section{Entusiasmo, desilusão e memória curta}
Os ciclos de entusiasmo e desilusão aparecem quando promessas, financiamento e capacidade observada deixam de andar juntos. A explicação não deve reduzir toda uma época a «as máquinas eram fracas» ou «as pessoas não perceberam». Importam dados disponíveis, custos, métodos, expectativas e aplicações que conseguiam pagar a tecnologia.

O leitor pode aplicar uma disciplina simples a qualquer anúncio. Identifica a tarefa que foi demonstrada, o recurso que a tornou possível e o obstáculo que continua por resolver. Se um vídeo mostra cinco sucessos e omite o número de tentativas, a taxa de sucesso é desconhecida. Se uma demonstração usa um operador escondido, mede o sistema humano e máquina em conjunto. São perguntas antigas que continuam a funcionar com nomes novos.

\begin{promptlabbox}{A manchete e o que foi realmente demonstrado}
\textbf{Objetivo e contexto:} analisar uma alegação pública sobre IA sem programar. Usa uma publicação com acesso ao artigo ou documentação original, e guarda a data.

\textbf{Procedimento:} escreve a manchete, a tarefa efetivamente avaliada, os dados, a ajuda disponível, a comparação e as limitações. Reescreve a manchete com o alcance que a evidência permite. Se não há denominador ou método, marca-o como desconhecido.

\textbf{Resultado esperado e validação:} outra pessoa deve conseguir ligar cada afirmação à fonte. «Não demonstrado» e «demonstrado falso» ficam separados.

\textbf{Falhas possíveis:} usar o anúncio comercial como avaliação independente, chamar consciência à conversação ou descartar um resultado útil por não ser AGI.
\end{promptlabbox}

O resto desta obra segue esta regra: compreender o mecanismo, experimentar numa tarefa e verificar o alcance. Não precisamos de esperar por uma definição definitiva de inteligência para contar as meias corretamente.
